%----------------------------------------------------------------------------------------

% Formatvorlage für Arbeiten von Mathias Weigert (Grundlage der Datei und allg. Struktur
% von Micha Schöneberger)

%--------------------------------------VERSION-------------------------------------------

% Version V0.10 - Erstellt von Mathias Weigert am 08.01.2014
%									Basis war die LaTex Dokumente der ZHAW
% Version V0.80 - Erstellt von Mathias Weigert am 24.08.2015
%									Versionsprung, damit Hauptdokument und Configurationsfiles nun selbe
%									Version besitzen. Alle benötigten Files sind nun in Unterverzeichnisse
%									ausgelagert.

% ---------------------------------LADEN DER META DATEN-----------------------------------

% Informationen über das Dokument, wie z.B. Titel, Autor, Jahr etc. werden in der
% Datei _Meta.tex definiert und können danach global verwendet werden.

%----------------------------------------------------------------------------------------

% Formatvorlage für Arbeiten von Mathias Weigert (Grundlage der Datei und allg. Struktur
% von Micha Schöneberger)

%--------------------------------------VERSION-------------------------------------------

% Version V0.75 - Erstellt von Mathias Weigert am 16.09.2013
%				  Bisher ungenützten "Balast" entfernt und zweite Matrikelnummer
%				  eingefügt.

%----------------------------------- INFORMATIONEN ---------------------------------------

% 	Definition von globalen Parametern, die im gesamten Dokument verwendet
% 	werden können (z.B auf dem Deckblatt etc.).
%-----------------------------------------------------------------------------------------

\newcommand{\titel}{Collector}	            % Titel des Dokumentes
\newcommand{\untertitel}{Semesterarbeit}	% Untertitel des Dokumentes
\newcommand{\autorFirst}{Mathias Weigert}	% Autor
\newcommand{\jahr}{2015}					% Jahr der Arbeit

\newcommand{\version}{Kick Off Semesterarbeit}				% Dokument Version
\newcommand{\gueltig}{2015}				    % Dokument Gültigkeit
\newcommand{\initial}{mwe}					% Initial des Autors
\newcommand{\freigabe}{}					% Freigegeben durch

\newcommand{\path}{D:\textbackslash Data\textbackslash Studium\textbackslash WorkSpace\textbackslash Collector\textbackslash Protokoll\textbackslash}	% Pfad des Files
\newcommand{\file}{KickOff.pdf}    	       % Name des Files

%-----------------------------------------------------------------------------------------

% ----------------------------------------------------------------------------------------

% ----------------------------------DOKUMENTENKOPF----------------------------------------

% Diese Vorlage basiert auf "scrreprt" aus dem koma-script.
% Die Option "draft" sollte beim fertigen Dokument ausgeschaltet werden.

\documentclass[
	9pt,				% Schriftgröße
	ngerman,			% für Umlaute, Silbentrennung etc.
	a4paper,       		% Papierformat
	oneside,			% einseitiges Dokument
	titlepage,			% es wird eine Titelseite verwendet
	final				% Status des Dokuments (final/draft)
]{scrbook}				% KOMA Script-Klassen: scrbook, scrreprt, scrartcl, scrlttr2.
						% Die Angabe einer Dokumentenklasse ist obligatorisch.

% ----------------------------------------------------------------------------------------

% -----------------------------PACKAGES EINLESEN------------------------------------------
% Weitere Packages, die benötigt werden, sind in die Datei Packages.tex "ausgelagert", um
% die Vorlage möglichst Übersichtlich zu halten.

%----------------------------------------------------------------------------------------

% Formatvorlage fuer Arbeiten von Mathias Weigert (Grundlage der Datei und allg. Struktur
% von Micha Schoeneberger)

%--------------------------------------VERSION-------------------------------------------

% Version V0.75 - Erstellt von Mathias Weigert am 16.09.2013
%				          Etwas aufgeraeumt, muss mal bei Zeit detailiert Ueberarbeitet werden.
%         V0.80 - Erstellt von Mathias Weigert am 24.08.2015
%                 Kleinere Anpassungen und nochmalieges "aufräumen" der einzelnen
%                 Packages.

%---------------------------------- INFORMATIONEN ---------------------------------------

% 	Definition von globalen Parametern, die im gesamten Dokument verwendet
% 	werden koennen (z.B auf dem Deckblatt etc.).
%----------------------------------------------------------------------------------------

\usepackage{cite}					% Verbesserte Zitierungsmoeglichkeiten
\usepackage{tocbasic}

%----------------------- FORMATIERUNG KOPF- UND FUSSZEILE -------------------------------

\usepackage
[automark,							% Automatische Kopfzeile
%headtopline,						% Linie ueber dem Seitenkopf
%plainheadtopline,					% Plain, Linie ueber dem Seitenkopf
headsepline,						% Linie zwischen Kopf und Textkoerper
ilines,								% Trennlinie linksbuendig ausrichten
%plainheadsepline,					% Plain, Linie zwischen Kopf und Textkoerper
footsepline,						% Linie zwischen Textkoerper und Fuss
plainfootsepline,   				% Plain, Linie zwischen Textkoerper und Fuss
%footbotline,						% Linie unter dem Fuss
%plainfootbotline   				% Plain, Linie unter dem Fuss
]{scrpage2}

%----------------------------------------------------------------------------------------


%-------------------------- ANPASSUNG LANDESSPRACHE--------------------------------------

\usepackage{ngerman}

%----------------------------------------------------------------------------------------


%--------------------------------- UMLAUTE-----------------------------------------------

\usepackage[T1]{fontenc}			% Dadurch ist auch sichergestellt, dass in einem PDF
									% Umlaute gefunden werden.
%\usepackage[utf8x]{luainputenc}
\usepackage[utf8]{inputenc}			% Unterstützung erweiterter Zeichensätze mit
									%  unterschiedlichen Kodierungen
									% (z. B. ae, oe,  usw.)
\usepackage{textcomp} 				% Text Companion Schriften (grosse Symbolbibliothek)

%----------------------------------------------------------------------------------------

\usepackage{pifont}					% PostScript standard Symbol and Dingbats fonts

%--------------------------------EIGENE REFERENZIERUNG-----------------------------------

\newcommand{\refTC}[1]{ \ref{#1} \nameref{#1}  }
% gibt aus z.B.: 1.1.3 Erwartetes Resultat (<Kapitelnummer> <Kapitelname>)

%----------------------------------------------------------------------------------------


%----------------------------------- GRAFIKEN--------------------------------------------

\usepackage[dvips,final]{graphicx}		% Einbinden von EPS-Grafiken [draft oder final]
\graphicspath{{pic/}} 					% Hier liegen die Bilder des Dokuments

%----------------------------------------------------------------------------------------

\usepackage{amsmath,amsfonts}			% Befehle aus AMSTeX fuer mathematische Symbole
										% z.B. \boldsymbol \mathbb
\usepackage{multirow}					% Packages fuer Tabellen

%-------------------------------- OWN COLORS---------------------------------------------

\usepackage{conf/rgbcolor}				% Eigene Farbdefinitionen nach RGB Chart
\usepackage[table]{xcolor}				% Um Tabellen einzuf�rben

%----------------------------------------------------------------------------------------

\usepackage{makeidx}					% Fuer Index-Ausgabe; \printindex

%------------------------------ QUELLCODE AUSGABE----------------------------------------

%************************************************************************************
%*                                                    Definitionen zur Quellcode Darstellung                                                   *
%************************************************************************************

\usepackage{listings}

\lstset {language=C++,
	keywordstyle={\color{Blue}\bfseries},
	commentstyle={\color{ForestGreen}\slshape},
	stringstyle={\color{Purple}},
	backgroundcolor={\color{PowderBlue}},
	showstringspaces=fales,
	stepnumber=2,
	numbers=left,
	numberstyle=\tiny}

\lstset {language=SQL,
	keywordstyle={\color{IndianRed}\bfseries},
	commentstyle={\color{ForestGreen}\slshape},
	identifierstyle=\ttfamily\color{CadetBlue}\bfseries,
	stringstyle={\color{Purple}},
	backgroundcolor={\color{Tan}},
	showstringspaces=fales,
	stepnumber=2,
	numbers=left,
	numberstyle=\tiny}						% Eigene Ausgabeformatierungen fuer Quellcode

%----------------------------------------------------------------------------------------


%-------------------------- ZEILENABSTAENDE UND SEITENRAENDER------------------------------

\usepackage{setspace}					% hat zur Benützung 3 Optionen: [singlespacing,
										% onehalfspacing, doublespacing]
\usepackage{geometry}
\usepackage{lineno}						% ermoeglicht das Numerieren der Zeilen

\usepackage{titlesec}
\titleformat{\chapter}
{\normalfont\huge\bfseries}{\thechapter.}{11pt}{\Huge}
\titlespacing*{\chapter}{0pt}{-40pt}{15pt}


% Symbolverzeichnis --------------------------------------------------------
% Symbolverzeichnisse bequem erstellen, beruht auf MakeIndex.
% makeindex.exe %Name%.nlo -s nomencl.ist -o %Name%.nls
% erzeugt dann das Verzeichnis. Dieser Befehl kann z.B. im TeXnicCenter
% als Postprozessor eingetragen werden, damit er nicht staendig manuell
% ausgefuehrt werden muss.
% Die Definitionen sind ausgegliedert in die Datei Abkuerzungen.tex.
% --------------------------------------------------------------------------
\usepackage[intoc]{nomencl}
  \let\abbrev\nomenclature
  \renewcommand{\nomname}{Abkürzungsverzeichnis}
  \setlength{\nomlabelwidth}{.25\hsize}
  \renewcommand{\nomlabel}[1]{#1 \dotfill}
  \setlength{\nomitemsep}{-\parsep}


%inserted november 2012 / Micha
\usepackage[normalem]{ulem}
\newcommand{\markup}[1]{\uline{#1}}


\usepackage{floatflt}							% Zum Umfliessen von Bildern

\usepackage[hyphens]{url}			% bricht lange URL's um, ohne dass der Link defekt ist
%\usepackage{url}

%-------------------------------PDF OPTIONEN----------------------------------------------

\usepackage[
bookmarks,
bookmarksopen=true,
pdftitle={\titel},
pdfauthor={\autorFirst},
pdfcreator={\autorFirst},
pdfsubject={\titel},
pdfkeywords={\titel},
colorlinks=true,
linkcolor=red, 							% einfache interne Verknuepfungen
anchorcolor=black,						% Ankertext
citecolor=blue, 						% Verweise auf Literaturverzeichniseintraege im Text
filecolor=magenta, 						% Verknuepfungen, die lokale Dateien oeffnen
menucolor=red, 							% Acrobat-Menuepunkte
urlcolor=cyan,
%linkcolor=black, 						% einffdrtext
%citecolor=black, 						% Verweise auf Literaturverzeichniseintraege im Text
%filecolor=black, 						% Verknuepfungen, die lokale Dateien öffnen
%menucolor=black,						% Acrobat-Menuepunkte
%urlcolor=black,
backref,
%pagebackref,
plainpages=false,						% zur korrekten Erstellung der Bookmarks
pdfpagelabels,							% zur korrekten Erstellung der Bookmarks
hypertexnames=false,					% zur korrekten Erstellung der Bookmarks
linktocpage 							% Seitenzahlen anstatt Text im Inhaltsverzeichnis verlinken
]{hyperref}

\usepackage[final]{pdfpages}			%Ermoeglicht andere PDF ins Dokument einzufuegen.

%-----------------------------------------------------------------------------------------

\usepackage{chngcntr}					% Zum fortlaufenden Durchnummerieren der Fussnoten

% fuer lange Tabellen
\usepackage{longtable}
\usepackage{array}
\usepackage{ragged2e}
\usepackage{lscape}

%Spaltendefinition rechtsbaendig mit definierter Breite
\newcolumntype{w}[1]{>{\raggedleft\hspace{0pt}}p{#1}}

% Formatierung von Listen aendern
\usepackage{paralist}
% \setdefaultleftmargin{2.5em}{2.2em}{1.87em}{1.7em}{1em}{1em}

\usepackage{lastpage}				% Dient zum ermitteln der Seiten eines Dokuments
	% lade ausgelagerte Packages

% ----------------------------------------------------------------------------------------

% -----------------------KOPF- UND FUSSZEILE, SEITENRÄNDER...-----------------------------
% Diese Einstellungen sind im Dokument Seitenstil ausgelagert.

% --------------------------------------------------------------------------ZEILENABSTAND-------------------------------------------------------------------------------------------

\onehalfspacing								% setzte den Zeilenabstand / Möglichkeiten: singlespacing, onehalfspacing, doublespacing

%------------------------------------------------------------------------------------------------------------------------------------------------------------------------------------------

			
% -------------------------------------------------------------------------SEITENRÄNDER--------------------------------------------------------------------------------------------

\geometry{									% definiert Papierformat und die Ränder
paper=a4paper,left=35mm,right=25mm,top=10mm, bottom=48mm
}

%------------------------------------------------------------------------------------------------------------------------------------------------------------------------------------------


% ------------------------------------------------------------------------KOPF- UND FUSSZEILE--------------------------------------------------------------------------------------

\pagestyle{scrheadings}							%
\renewcommand*{\chapterpagestyle}{scrheadings}			% Kopf- und Fußzeile auch auf Kapitelanfangsseiten


\renewcommand{\headfont}{\normalfont}				% definiert die Schriftart für die Kopfzeile (z.B: auch möglich: \sffamily)

% Kopfzeile ----------------------------------------------------------------
%\ihead{\large{\textsc{\titel}}\\[1ex] \textit{\headmark}}
\ihead{\textit{\headmark}}
\chead{\qquad \qquad \qquad \qquad \qquad \qquad \qquad \quad \qquad \qquad \qquad \qquad \qquad \qquad \qquad   \includegraphics[scale=0.08]{pic/ZHAW_logo.png}}%\ohead{\includegraphics[scale=0.08]{ZHAW_logo.png}} % setze Aussenseite der Kopfzeile
\setlength{\headheight}{30mm} 						% Höhe der Kopfzeile
\setheadwidth[0pt]{textwithmarginpar} 				% Kopfzeile über den Text hinaus verbreitern
\setheadsepline[text]{0.4pt} 						% Trennlinie unter Kopfzeile

\ifoot{\autorFirst}								% setze Innenseite der Fusszeile	
\cfoot{}									% setze Center-Fusszeile
\ofoot{\pagemark}								% setze Aussenseite Fusszeile
\setfootsepline[text]{0.4pt} 						% Trennlinie über Fusszeile

%------------------------------------------------------------------------------------------------------------------------------------------------------------------------------------------
%  INFO:

%  @\pagemark								% Kapitelname
%  @ihead									% Mit dieser Angabe wird die Innenseite ("i" = "inside") des Seitenkopfs definiert
%  @chead									% Mit dieser Angabe wird die Mitte ("c" = "center") des Seitenkopfs definiert
%  @ohead									% Mit dieser Angabe wird die Außenseite ("o" = "outside") des Seitenkopfs definiert
%  @ifoot									% Mit dieser Angabe wird die Innenseite des Fusszeile definiert
%  @cfoot									% Mit dieser Angabe wird die Mitte des Fusszeile definiert
%  @ofoot									% Mit dieser Angabe wird die Aussenseite des Fusszeile definiert

% siehe auch: http://www.jochen-lipps.de/latex/kopfzeile.php
% siehe auch: http://www2.informatik.hu-berlin.de/~piefel/LaTeX-PS/Archive-2004/V07-footnote.pdf
%------------------------------------------------------------------------------------------------------------------------------------------------------------------------------------------


% --------------------------------------------------------SCHUSTERJUNGEN UND HURENKINDER----------------------------------------------------------------------------------

\clubpenalty = 10000							% Disable single lines at the start of a paragraph (Schusterjungen)
\widowpenalty = 10000 							% Disable single lines at the end of a paragraph (Hurenkinder)
\displaywidowpenalty = 10000						% Disable single lines at the end of a paragraph (Hurenkinder)

%------------------------------------------------------------------------------------------------------------------------------------------------------------------------------------------

% -----------------------------------------------------------AUSGABE QUELLCODE FORMATIEREN---------------------------------------------------------------------------------

\lstset{numbers=left, numberstyle=\tiny, numbersep=5pt, breaklines=true}
\lstset{emph={square}, emphstyle=\color{red}, emph={[2]root,base}, emphstyle={[2]\color{blue}}}									%------------------------------------------------------------------------------------------------------------------------------------------------------------------------------------------

\frenchspacing								% erzeugt ein wenig mehr Platz hinter einem Punkt
\counterwithout{footnote}{chapter}					% Fußnoten fortlaufend durchnummerieren
	% lade spezifische Kopf- und Fusszeilen

% ----------------------------------------------------------------------------------------

\begin{document}

\setlength{\parindent}{0pt}		% sorgt dafür das es keinen Einzug am Absatzanfang gibt
\setcounter{secnumdepth}{3}		% legt die Tiefe bei Nummerierung von Überschriften fest
\setcounter{tocdepth}{3}		% legt Schachtelungstiefe bei Inhaltsverzeichnis fest

% ----------------------------------------------------------------------------------------

% ---------------------DECKBLATT, INHALT & VERSIONSVERWALTUNG-----------------------------

%%----------------------------------------------------------------------------------------

% Formatvorlage für Arbeiten von Mathias Weigert (Grundlage der Datei und allg. Struktur
% von Micha Schöneberger)

%--------------------------------------VERSION-------------------------------------------

% Version V0.75 - Erstellt von Mathias Weigert am 16.09.2013
%				  				Der besseren Übersicht wegen Informationen aus zwei Seiten verteilt
% 				V0.80 - Erstellt von Mathias Weigert am 24.08.2015
%									Deckblatt auf einer Seite mit etwas weniger Information. Die Zukunft
%									wird so aussehen, dass ich verschiedene Deckblätter erstelle welche
%									nach Bedarf ins Dokument eingebunden werden.

%------------------------------------SEITENLAYOUT-----------------------------------------

% Informationen Über das Dokument, wie z.B. Titel, Autor, Matrikelnr. etc werden in der
% Datei _Meta.tex definiert und können danach global verwendet werden.

%\thispagestyle{empty} 		% Dokument ohne Kopfzeile und ohne Seitennummerierung.
\thispagestyle{plain} 		% Mit Seitennummerierung (Standard)
%\thispagestyle{headings} 	% Generiert eine automatische Kopfzeile mit Seitenzahl und
							% Zwischenüberschriften
%\thispagestyle{myheadings} % Erlaubt die Erstellung eigener Kopf- und Fußzeilen
%\thispagestyle{fancy} 		% Erlaubt die Verwendung der in dem Paket "fancyhdr"
							% definierten Befehle zur Erstellung eigener Kopf- und
							% Fußzeilen

%-----------------------------------------------------------------------------------------

%----------------------------------SETZE TITELBLATT---------------------------------------

\begin{titlepage}

\begin{center}
\includegraphics[scale=0.25]{pic/ZHAW_Titel}			% binde die Grafik ein
\\[12ex]
\huge{\textbf{\textsc{\titel}}}							% setze Schriftgrösse für Titel
\\[1.5ex]
\LARGE{\textbf{\untertitel}}							% setze Schriftgrösse für Untertitel
\\[6ex]
\normalsize												% setze Schriftgrösse zurück auf Normalschrift

\begin{tabular}{w{3cm}p{6cm}}
\\[12ex]Autor: & \quad \autorFirst						% Angabe des Studienbereichs
\end{tabular}
\\[24ex]
\copyright\ \jahr\\
\end{center}

\singlespacing											% einfacher Zeilenabstand
\small													% setze kleine Schrift
\noindent												%	Sorgt am Anfang eines Absatzes dafür, daß die
														% erste Zeile nicht eingerückt wird.


%----------------------------GNU FREE DOCUMENTATION LICENSE---------------------------------
Es wird die Erlaubnis gegeben dieses Dokument zu kopieren, verteilen und/oder zu verändern
unter den Bedingungen der GNU Free Documentation License, Version 1.3 oder einer späteren, von
der Free Software Foundation veröffentlichten Version; mit den Unveränderlichen Abschnitten
DEREN TITEL	AUFGEZÄHLT sind, mit den Vorderseitentexten die AUFGEZÄHLT sind, und mit den
Rückseitentexten die AUFGEZÄHLT sind. Eine Kopie dieser Lizenz ist in	dem Abschnitt enthalten,
der mit \href{http://www.gnu.org/copyleft/fdl.html}{GNU Free Documentation License} betitelt ist.

%--------------------------------------COPYRIGHT-------------------------------------------
%Dieses Werk einschließlich seiner Teile ist \textbf{urheberrechtlich geschützt}. Jede
%Verwertung außerhalb der engen Grenzen des Urheberrechtsgesetzes ist ohne Zustimmung des
%Autors unzulässig und strafbar. Das gilt insbesondere für Vervielfältigungen, Übersetzungen,
%Mikroverfilmungen sowie die Einspeicherung und Verarbeitung in elektronischen Systemen.

\end{titlepage}
		% Deckblatt

%\tableofcontents				% Inhaltsverzeichnis

%\chapter*{Version}

Angabe der aktuellen Version des Dokuments. Alle Änderungen, werden mit Versionsnummer, Datum, Autor und Informationen betreffend der Änderung dokumentiert.\\

\begin{tabular}{|c|l|c|l|}
	\rowcolor{black} {\color{white}\textbf{Version}} & {\color{white}\textbf{Datum}} & {\color{white}\textbf{Autor}} & {\color{white}\textbf{Bemerkung}} \\
	0.01 & 30.11.15 & M. Weigert & Initiales Dokument \\ \hline
	\rowcolor{DarkSeaGreen} 0.02 & 30.11.15 & M. Weigert & Aufgabe hinzugefügt \\ \hline
	0.03 & 02.12.15 & M. Weigert & Neue LaTex Konfiguration, Literaturverzeichnis, erste Recherche \\ \hline
	\rowcolor{DarkSeaGreen} 0.04 & 04.12.15 & M. Weigert & Kapitelstruktur angelegt. Recherche fortgeführt \\ \hline
	0.04 & 03.02.16 & M. Weigert & Zeitmanagement und Feinstruktur \\ \hline
	\rowcolor{DarkSeaGreen} 0.05 & 06.02.16 & M. Weigert & Neue Datenstruktur und Anforderungen dokumentiert. \\ \hline
\end{tabular}

\chapter*{Zeitmanagement}

\begin{tabular}{|c|l|c|c|l|}
	\rowcolor{black} {\color{white}\textbf{Datum}} & {\color{white}\textbf{Start}} & {\color{white}\textbf{Ende}} & {\color{white}\textbf{Zeit (min.)}} & {\color{white}\textbf{Bemerkung}} \\
	03.02.16 & 20:30 & 21:00 & 30 & Zeitmanagement und Feinplanung (Basis) \\ \hline
	\rowcolor{DarkSeaGreen} 06.02.16 & 11:00 & 12:45 & 105 & Angefangen mit der Anforderungsanalyse \\ \hline
	06.02.16 & 13:30 & & & Weiter an Anforderungsanalyse gearbeitet. \\ \hline
\end{tabular}	% Versionsverwaltung

% ----------------------------------------------------------------------------------------

% -----------------------------INHALT DES DOKUMENTS---------------------------------------

% Hier können jetzt die einzelnen Kapitel implementiert werden. Sie müssen in den
% entsprechenden .TEX-Dateien vorliegen. Die Dateinamen können natürlich angepasst werden.

% ----------------------------------------------------------------------------------------

\begin{tabular}{ll}
	\textbf{Protokoll:} & Design Review Seminararbeit \\
	\textbf{Anwesend:} & R. Knaack, M. Reiser, M. Weigert \\
	\textbf{Protokollführer:} & M. Weigert \\
\end{tabular}

\begin{enumerate}
	\item Ist der Auftrag für die Semesterarbeit von der Studentin oder dem Studenten korrekt erfasst? \\
	{\color{DarkSlateBlue}LOREM IPSUM DOLOR}
	\item Sind die Ausgangslage und das Umfeld ausreichend analysiert, berücksichtigt und bearbeitet? \\
	{\color{DarkSlateBlue}LOREM IPSUM DOLOR}
	\item Wurde eine systematische Recherche durchgeführt?\\
	{\color{DarkSlateBlue}LOREM IPSUM DOLOR}
	\item Wurde ein klares Konzept erarbeiotet und klar dargestellt?\\
	{\color{DarkSlateBlue}LOREM IPSUM DOLOR}
	\item Wurden alternative Lösungen betrachtet? \\
	{\color{DarkSlateBlue}LOREM IPSUM DOLOR}
	\item Entsprechen das Arbeits- und Lösungskonzept den Anforderungen an eine Semesterarbeit? \\
	{\color{DarkSlateBlue}LOREM IPSUM DOLOR}
	\item Ist die Lösung grundsätzlich für die Auftraggeberin bzw. den Auftraggeber akzeptabel? \\
	{\color{DarkSlateBlue}LOREM IPSUM DOLOR}
	\item Ist das Konzept bzw. die Lösung technisch und terminlich im Rahmen der verbleibenden Zeit umsetzbar? \\
	{\color{DarkSlateBlue}LOREM IPSUM DOLOR}
	\item Sind die nächsten Arbeitsschritte klar formuliert?\\
	{\color{DarkSlateBlue}LOREM IPSUM DOLOR}
\end{enumerate}

\textbf{M. Reiser:}
\begin{itemize}
	\item INSERT ITEM HERE
\end{itemize}

\textbf{R. Knaack}
\begin{itemize}
	\item INSERT ITEM HERE
\end{itemize}

% --------------------------------------ANHANG--------------------------------------------

% Die Inhalte des Anhangs werden analog zu den Kapiteln implementiert; über die Datei
% Anhang.tex

%\begin{appendix}
%	\clearpage
%	\pagenumbering{Roman}
%	\renewcommand{\thesection}{\arabic{section}}
%	\chapter{Anhang}

\section{Konventionen}

Für die Beschreibung von Eingaben oder Beschriftungen in der App wird folgendes Format benutzt. \\
{\color{IndianRed}\texttt{Beispiel einer Eingabe durch den User oder Beschriftung in der App.}}

\section{Verwendete Tools \& Software}

\subsection{Dokumentation \& Präsentation}

Alle Dokumente und Präsentationen wurden in \LaTeX \ geschrieben. Dazu wurde \href{http://texstudio.sourceforge.net}{TeXstudio} in der Version 2.10.4 genutzt. 

\subsection{Programmierung}

Die App wurde im \href{https://developer.android.com/studio/index.html}{Android Studio} 1.5.1 von Google erstellt.\\
Zusätzlich wurde für den CSV Export die Library \href{http://opencsv.sourceforge.net/}{Opencsv} in der Version 3.7 genutzt.

\subsection{Code und Versionsverwaltung}

Für die Versionsverwaltung aller Dokumente und des Programmcode wurde auf GitHub ein \href{https://github.com/MWeigert/Collector}{Repository} eingerichtet.

\subsection{Datenbank}

Für die Erstellung des Datenbank Design wurde der \href{http://fabforce.net/dbdesigner4/}{DB Designer 4} von fabFORCE.net verwendet.

\subsection{GUI}

Für die Funktion und das Design der GUI wurde das Webtool \href{https://www.fluidui.com}{Fluid} genutzt.

\subsection{UML}

Sämtliche UML wurden mit dem freien Tool \href{http://www.umlet.com}{Umlet} erstellt.  

\section{Design Entscheidungen}

\begin{tabular}{|c|l|c|l|}
	\rowcolor{black} {\color{white}\textbf{ID}} & {\color{white}\textbf{Datum}} & {\color{white}\textbf{Autor}} & {\color{white}\textbf{Bemerkung}} \\
	001 & 12.03.16 & M. Weigert & Hauptsprache für App wird Englisch \\ \hline
	\rowcolor{DarkSeaGreen} 002 & 12.03.16 & M. Weigert & Es wird auf jeden Zurück/Back Button verzichtet. \\ \hline
	003 & 25.03.16 & M. Weigert & Extra Tabelle Parental um internationale Jugendfreigaben zu managen. \\ \hline
	\rowcolor{DarkSeaGreen} 004 & 09.04.16 & M. Weigert & Verzicht auf die Tabelle Parental um die App nicht unötig zu komplizieren. \\ \hline
\end{tabular}
%	\bibliography{bib/mybib.bib}
%	\bibliographystyle{unsrt}
%\end{appendix}

% ----------------------------------------------------------------------------------------

\end{document}
